% SOURCE_AUTHORITY: sga5_fr_workpass.tex, lines 9107--9169 (LF-normalized unit).
% SOURCE_SHA256: 791F4EFFC5E02832D5D77ED03518C8156D6F07E4C8238B03545DB93D883FBB28
\begin{statement}{Corolário 1.2}
Seja $L$ um $A_S$-módulo. Para todo inteiro $i$, o morfismo de imagem inversa
\[
H^i(S,L)\longrightarrow H^i(X,p^*L)
\]
é um isomorfismo.
\end{statement}

\emph{Demonstração.} Por (1.1), o morfismo de adjunção $L\to \R p_*(p^*L)$ é um isomorfismo. Como $H^*(X,p^*L)\simeq H^*(S,\R p_*p^*(L))$, segue-se a afirmação.

\begin{statement}{Proposição 1.3}
Suponhamos que $E$ tem posto constante $r$. Para todo $A_S$-módulo $L$, verificam-se as seguintes propriedades.
\begin{enumerate}[label=(\roman*)]
\item
\begin{enumerate}[label=\alph*)]
\item
\[
\R^i q_*(q^*L)=0\quad\text{se }i\ne 0,2r-1.
\]
\item O morfismo de adjunção
\[
L\longrightarrow q_*q^*(L)
\]
é um isomorfismo.
\item Denotando por $Y$ O complementar de $U$ em $X$, o morfismo
\[
\R^{2r-1}q_*(q^*L)\longrightarrow L\otimes\mu_S^{\otimes -r}
\]
Composto do morfismo de borda
\[
\R^{2r-1}q_*(q^*L)\longrightarrow \R_Y^{2r}p_*(p^*L)
\tag{1}
\]
e isomorfismo de pureza (SGA 4 XVIII)
\[
\R_Y^{2r}p_*(p^*L)\longrightarrow L\otimes\mu_S^{\otimes -r}
\]
é um isomorfismo.
\end{enumerate}
\item
\begin{enumerate}[label=\alph*)]
\item
\[
\R_!^i q(q^*L)=0\quad\text{se }i\ne 1,2r.
\]
\item O morfismo de borda
\[
L\longrightarrow \R_!^1 q(q^*L)
\]
proveniente da longa sequência exata, onde $f:Y\to S$ é a canónica projeção,
\[
\cdots\longrightarrow \R_!^i q(q^*L)\longrightarrow \R_!^i p(p^*L)
\longrightarrow \R^i f_*(f^*L)\longrightarrow \R_!^{i+1}q(q^*L)\longrightarrow\cdots
\]
é um isomorfismo.
\item O morfismo de traço
\[
\R_!^{2r}q(q^*L)\longrightarrow L\otimes\mu_S^{\otimes -r}
\]
é um isomorfismo.
\end{enumerate}
\end{enumerate}
\end{statement}
